%\documentclass[twocolumn]{article}\
\documentclass{article}
\usepackage{amsmath, amssymb, cancel, mathtools, bm}
\usepackage[left=.5in, right=.5in, top=1in, bottom=1in]{geometry}
\usepackage[most]{tcolorbox}
\usepackage[skip=1em,indent=0pt]{parskip}
\setlength{\parindent}{0pt}
\newcommand{\svec}[1]{\underset{\sim}{\bm{#1}}} % Vector symbol (tilde under vector)
\DeclareMathOperator{\EX}{\mathbb{E}} % Expected Value symbol
\DeclareMathOperator{\ext}{\EX_{\theta}} % Expected Value symbol
\DeclareMathOperator{\vart}{Var_{\theta}} % Expected Value symbol
\newcommand{\indep}{\perp\!\!\!\!\perp} % Independence symbol
\newcommand{\real}{\mathbb{R}} % Simplified 'Reals' indicator
\newcommand{\pto}{\overset{P}{\to}}
\newcommand{\asto}{\overset{a.s.}{\to}}
\newcommand{\rthto}{\overset{L^r}{\to}}
\newcommand{\dto}{\overset{D}{\to}}
\newcommand{\xtb}{x^{\top}\svec{\beta}}
\newcommand{\xitb}{x_i^{\top}\svec{\beta}}
% Force all aggregate symbols to always put values above/below
\let\Oldint=\int
\let\Oldsum=\sum
\let\Oldprod=\prod
\let\Oldbigcup=\bigcup
\let\Oldbigcap=\bigcap
\let\Oldlim=\lim
\renewcommand{\int}{\Oldint\limits} 
\renewcommand{\sum}{\Oldsum\limits} 
\renewcommand{\prod}{\Oldprod\limits} 
\renewcommand{\bigcup}{\Oldbigcup\limits} 
\renewcommand{\bigcap}{\Oldbigcap\limits} 
\renewcommand{\lim}{\Oldlim\limits} 
\newcommand{\infint}{\int_{-\infty}^{\infty}}
\newcommand{\iiddist}{\overset{\mathrm{iid}}{\sim}}
\newcommand{\norm}[1]{\left\lVert#1\right\rVert}
\newcommand{\boldred}[1]{\textbf{\textcolor{red}{#1}}}

\author{RJ Cass}
\title{STAT 537 - HW 3}

\begin{document}

\maketitle

\section{MLE for $\Sigma$}

Goal: Find the MLE of $\Sigma$ in the 2x2 case.

Likelihood: $L = (2\pi)^{-n/2}|\Sigma|^{-n/2}e^{-\frac{1}{2}(y-x\beta)^{\top}\Sigma^{-1}(y-x\beta)}$

Log-likelihood: $l = -\frac{n}{2}ln(2\pi) - \frac{n}{2}ln(|\Sigma|) - \frac{1}{2}(y-x\beta)^{\top}\Sigma^{-1}(y-x\beta)$

Ok, at this point I needed to use chat to understand how to take the derivative of the last portion. Basically, we must first write it as a trace (which we can do beacuse the result of this function is a scalar, and for a scalar $a, a = tr(a)$):

$\sum[(y-x\beta)\Sigma^{-1}(y-x\beta)] = tr((y-x\beta)\Sigma^{-1}(y-x\beta))$

Due to the properties of the trace (apparently called the cyclis property), we can rearrange the matrices within the trace:

$tr(\sum[(y-x\beta)\Sigma^{-1}(y-x\beta))] = tr(\Sigma^{-1}(y-x\beta)(y-x\beta)^{\top})$

Now we use the propert of the trace in matrix calculus where $\frac{\partial}{\partial A} tr(AB) = B^{\top}$:

$\frac{\partial}{\partial \Sigma^{-1}} tr(\Sigma^{-1}(y-x\beta)(y-x\beta)^{\top}) = ((y-x\beta)(y-x\beta)^{\top})^{\top} = (y-x\beta)(y-x\beta)^{\top}$

I also needed chat to learn that $\frac{d}{d \Sigma^{-1}} ln(|\Sigma|) = -\Sigma$, and $ln(|\Sigma|) = -ln(|\Sigma^{-1}|)$

$\frac{\partial}{\partial \Sigma^{-1}} -\frac{n}{2}ln(2\pi) + \frac{1}{2}ln(|\Sigma^{-1}|) - \frac{1}{2}(y-x\beta)^{\top}\Sigma^{-1}(y-x\beta) = 0 - \frac{n}{2}\Sigma + \frac{1}{2}(y-x\beta)(y-x\beta)^{\top}$

Setting this equal to 0:

$\frac{n}{2}\Sigma = \frac{1}{2}(y-x\beta)(y-x\beta)^{\top}$

So the MLE for $\Sigma$ is:

$\hat{\Sigma} = \frac{(y-x\beta)(y-x\beta)^{\top}}{n}$

\section{MLEs of $\sigma^2_I$ and $\sigma^2_e$}

$(y-x\beta)(y-x\beta)^{\top} = 
    \begin{pmatrix} (y_1 - x\beta_1)^{\top} \\ (y_2 - x\beta_2)^{\top} \end{pmatrix}
    \begin{pmatrix} y_1 - x\beta_1 & y_2 - x\beta_2 \end{pmatrix}
    = 
    \begin{pmatrix}
        (y_1 - x\beta_1)^{\top}(y_1 - x\beta_1) & (y_1 - x\beta_1)^{\top}(y_2 - x\beta_2) \\
        (y_2 - x\beta_2)^{\top}(y_1 - x\beta_1) & (y_2 - x\beta_2)^{\top}(y_2 - x\beta_2)
    \end{pmatrix}
$

From notes in class we know: $\Sigma = \begin{pmatrix} 
    \sigma_e^2 + \sigma_I^2 & \sigma_I^2 \\ 
    \sigma_I^2 & \sigma_e^2 + \sigma_I^2 
\end{pmatrix}$

So, we can set them equal to each other: 

$
\begin{pmatrix} 
    \sigma_e^2 + \sigma_I^2 & \sigma_I^2 \\ 
    \sigma_I^2 & \sigma_e^2 + \sigma_I^2 
\end{pmatrix} = 
\frac{1}{n}\begin{pmatrix}
    (y_1 - x\beta_1)^{\top}(y_1 - x\beta_1) & (y_1 - x\beta_1)^{\top}(y_2 - x\beta_2) \\
    (y_2 - x\beta_2)^{\top}(y_1 - x\beta_1) & (y_2 - x\beta_2)^{\top}(y_2 - x\beta_2)
\end{pmatrix}
$

The off-diagonal terms are the same: $(y_2 - x\beta_2)^{\top}(y_1 - x\beta_1) = ((y_1 - x\beta_1)^{\top}(y_2 - x\beta_2))^{\top}$, so we get: 

$\hat{\sigma}^2_I = \frac{1}{n}(y_1 - x\beta_1)^{\top}(y_2 - x\beta_2)$

However, the diagonal terms of our ingot covariance are the same, and the sampled data likely are not going to give the same values there. So, we use the average of those 2 values:

$\hat{\sigma}_e^2 + \hat{\sigma}_I^2 = \frac{1}{n}\frac{(y_1 - x\beta_1)^{\top}(y_1 - x\beta_1) + (y_2 - x\beta_2)^{\top}(y_2 - x\beta_2)}{2}$

Thus, we have the following MLEs:

$\hat{\sigma}^2_I =\frac{1}{n}(y_1 - x\beta_1)^{\top}(y_2 - x\beta_2)$

$\hat{\sigma}_e^2 = \frac{1}{n}\frac{(y_1 - x\beta_1)^{\top}(y_1 - x\beta_1) + (y_2 - x\beta_2)^{\top}(y_2 - x\beta_2)}{2} - \frac{1}{n}(y_1 - x\beta_1)^{\top}(y_2 - x\beta_2)$

\end{document}